import seaborn as sns
import matplotlib.pyplot as plt

# Korrelationsmatrix berechnen
corr_matrix = df_california.corr()

# Korrelationsplot erstellen
plt.figure(figsize=(10, 8))
sns.heatmap(
    corr_matrix,
    annot=True,           # Werte in den Zellen anzeigen
    cmap='coolwarm',      # Farbpalette: blau-weiß-rot
    vmin=-1, vmax=1,      # Wertebereich -1 bis 1
    square=True,          # Quadratische Zellen
    fmt='.2f'             # 2 Nachkommastellen
)

plt.title('Korrelationsplot des California-Housing-Datensatzes', fontsize=14)
plt.tight_layout()
plt.show()

# Korrelationen mit der Zielvariable extrahieren
corr_with_target = corr_matrix['MedHouseVal'].sort_values(ascending=False)

print("\nKorrelationen mit der Zielvariable MedHouseVal:")
print(corr_with_target)

# Variable mit der größten Korrelation
highest_corr = corr_with_target.index[1]  # Index 0 = MedHouseVal (1.0)
highest_corr_value = corr_with_target.iloc[1]

print(f"\nDie größte Korrelation mit MedHouseVal zeigt '{highest_corr}' "
      f"mit einem Wert von {highest_corr_value:.4f}")